\documentclass[italian,12pt]{article}

\def\Title{Verbale interno}
\def\Subtitle{Riunione interna settimanale}
\def\Author{7Last}
\def\Date{2024-05-22}
\def\Version{v1.0}

%pacchetti extra da scaricare dblfloatfix, fancyhdr
\usepackage[left=2cm, right=2cm, bottom=3cm, top=3cm]{geometry}
\usepackage{fancyhdr}%creazione header-footer
\usepackage{graphicx} %serve per inserire immagini
\graphicspath{ {../../../logo/} }
%\usepackage{dblfloatfix} %serve per posizionare gli elementi dove si vuole
\usepackage[hidelinks]{hyperref} %serve per i link
\usepackage{tikz}
\usepackage{tgadventor} % font
\usepackage[useregional=numeric,showseconds=true,showzone=false]{datetime2}
\usepackage{caption}

\usepackage{hyperref}
\usepackage{tocloft}
\usepackage{titlesec}
\usepackage{color}
\usepackage{ulem}
\usepackage{pgfplots}
\usepackage{pgf-pie}
\usepackage[italian]{babel}
\usepackage{comment}
\usepackage{tabularx}
\usepackage{longtable}
\usepackage{float}
\usepackage{amsmath}

\linespread{1.2}
\captionsetup[table]{labelformat=empty}
\geometry{headsep=1.5cm}

\renewcommand{\contentsname}{Indice}
\renewcommand\familydefault{\sfdefault}

\begin{document}

\newgeometry{left=2cm,right=2cm,bottom=2.1cm,top=2.1cm}
\begin{titlepage}
    \vspace*{.5cm}

    \vspace{2cm}
    {
        \centering
        {\bfseries\huge \Title\par}
        \bigbreak
        {\bfseries\large \Author\par}
        \bigbreak
        {\Version\par}
        \vfill

        \begin{tikzpicture}[remember picture,overlay]

            \fill[blue!20!black] (current page.south west) -- ++(10cm,0) -- ++(-10cm,15cm) -- cycle;
            \fill[orange] (current page.south east) -- ++(-18cm,0) -- ++(21.6cm,22cm) -- cycle;

            \clip (0,-3cm) circle (2.5cm) node (current page.center) {\includegraphics[width=5cm]{logo.jpg}};
        \end{tikzpicture}

    }

    \vfill

\end{titlepage}

\restoregeometry






















\newpage

%---------------header------------------%
\pagestyle{fancy}
\fancyhead{} % pulizia degli header
\lhead{%
\begin{tikzpicture}
    \clip (0,0) circle (0.5cm);
    \node at (0,0) {\includegraphics[width=1cm]{logo}};
\end{tikzpicture}%
}
\chead{\vspace{\fill}\Title\vspace{\fill}}
\rhead{\vspace{\fill}\Version\vspace{\fill}}
%quad è una spaziatura
%---------------------------------------%

\begin{table}[!h]
	\caption{Versioni}
	\footnotesize
	\begin{center}
		\begin{tabular}{ l l l l l }
			\hline \\[-2ex]
			Ver. & Data       & Redattore          & Verificatore       & Descrizione  \\
			\\[-2ex] \hline \\[-1.5ex]
			1.0  & 23/04/2024 & Leonardo Baldo  & Antonio Benetazzo & Stesura verbale \\
			\\[-1.5ex] \hline
		\end{tabular}
	\end{center}
\end{table}

\newpage

\tableofcontents

\newpage

\section{Dettagli della riunione}


\textbf{Sede della riunione}: Google Meet\\
\textbf{Orario di inizio}: 15:30\\
\textbf{Orario di fine}: 16:30\\


\begin{flushleft}
	\begin{table}[!h]
	\begin{tabular}{ |l|l|l| } 
		\hline
		\textbf{Partecipante} & \textbf{Ruolo}       & \textbf{Presenza} \\
		\hline 
		Antonio Benetazzo     & Verificatore         & Presente          \\
		Davide Malgarise      & Amministratore       & Presente          \\
		Elena Ferro           &                      & Presente          \\
		Leonardo Baldo        & Redattore	         & Presente          \\
		Matteo Tiozzo         &                      & Assente           \\
		Raul Seganfreddo      &                      & Presente          \\
		Valerio Occhinegro    &                      & Presente          \\
		\hline
	\end{tabular}
	\end{table}
	\textbf{Partecipanti esterni}: Daniele Zorzi.
\end{flushleft}
\section{Ordine del giorno}
\begin{itemize}
	\item Lavoro svolto
	\item Apache Kafka vs Redpanda
	\item Obiettivi secondo sprint
	\item Decisioni prese e conclusioni
\end{itemize}
\newpage
\section{Verbale}
\subsection{Lavoro svolto}
L'incontro inizia con una spiegazione del codice sviluppato fino al presente, seguito da una visualizzazione dei risultati ottenuti.
L'azienda considera il lavoro svolto finora ben fatto. In particolare, è stata apprezzata l'integrazione con Kafka e lo sviluppo del simulatore per un sensore aggiuntivo, entrambi ritenuti opzionali per questo sprint.
\subsection{Apache Kafka vs Redpanda}
Durante l'integrazione del codice simulatore di sensori con Kafka, è emersa un'ulteriore alternativa, non ancora valutata, rispetto \textit{Apache Kafka}: \textbf{\textit{Redpanda}}. Dopo un'attenta comparazione e analisi dei due strumenti, il team ha deciso di proporre all'azienda l'utilizzo di quest'ultimo. La prononente richiede di produrre un documento che illustri i vantaggi e gli svantaggi di entrambi gli strumenti, per poter valutare le due opzioni. Le conclusioni verranno discusse durante il prossimo SAL. 
\subsection{Obiettivi secondo sprint}
Per il secondo sprint sono stati definiti i seguenti obiettivi:
\begin{itemize}
	\item continuazione programma \textit{Python} per la generazione dei dati di test;
	\item aggiungere alcune tabelle di aggregazione dati su \textit{ClickHouse};
	\item visualizzazione su \textit{Grafana} di almeno un sensore;
	\item continuazione della documentazione;
	\item stesura \textit{Analisi dei Requisiti}.
\end{itemize}
\newpage
\subsection{Decisioni prese e conclusioni}
Discutiamo della durata degli sprint, concordando con l'azienda di mantenere una durata di 2 settimane, riservandoci la possibilità di modificarla in futuro.
La riunione si conclude dandoci appuntamento per il prossimo SAL, fissato in data 2024-05-06.
\begin{table}[b]
	\begin{tabular}{@{}p{.5in}p{4in}@{}}
		Data:  & \hrulefill \\
			   &     		\\
			   &     		\\
		Firma: & \hrulefill \\
	\end{tabular}
	\end{table}

\end{document}
